% Post-CP21 candidate. Exact universal b^4 minor retained.
\subsection{Degree-one full image on every nonordinary principal surface}
\status{Proved and independently checked.}

\begin{lemma}[The two-rank-one polar tensor]
\label{lem:p-rank-one-genus-two-polar-stable-candidate}
Let $C$ be a smooth genus-two curve over an algebraically closed field
of characteristic two. If its Jacobian has two-rank one, the alternating
polar three-form of its Coble cubic is stable.
\end{lemma}
\begin{proof}
Begin with the degree-five Weierstrass equation and the triangular
Cartier matrix calculated in
Lemma~\ref{lem:p-rank-zero-genus-two-polar-stable-candidate}.
Stable rank one means exactly one of $h_1,h_2$ is nonzero. If $h_2=0$,
$h$ is already linear. If $h_1=0$ and $h_2\ne0$, translate the double
root of $h$ to zero, so that $h=h_2x^2$. Translate $y$ by a constant
to make $f_5(0)=0$. Smoothness at $(0,0)$ implies $f'_5(0)\ne0$,
because both the $y$-derivative and $h'(0)$ vanish there. Now set
\[
                         z=x^{-1},\qquad v=yx^{-3}.
\]
The transformed equation is
\[
                  v^2+h_2zv=z^6f_5(z^{-1}).
\]
Its right side has degree exactly five: the original constant
coefficient is zero and its linear coefficient is nonzero. Thus
its $h$ is linear, and the new infinity is the other Weierstrass point.
This explicit transformation needs no assertion about the number of
points in an inseparable branch scheme.

In the linear case, translate its root to zero, and translate $y$
by a constant to arrange $f_5(0)=0$. Nonzero rescalings of $x,y$
normalize $h=x$ and the leading coefficient of $f_5$ to one:
if the initial coefficients are $h_1x$ and $f_5^{\rm lead}x^5$,
one chooses the $x$ scale with cube $h_1^2/f_5^{\rm lead}$ and the
$y$ scale equal to $h_1$ times that $x$ scale.
The substitution $y\mapsto y+ux^2$ adds $u^2x^4+ux^3$, so removes
the fourth-degree term. The substitution $y\mapsto y+vx$ adds
$(v^2+v)x^2$, so removes the quadratic term. Both equations for
$u,v$ have solutions over the algebraically closed field.
Every curve in question consequently has the form
\begin{equation}\label{eq:p-rank-one-normal-form-candidate}
                   y^2+xy=x^5+a x^3+b x,\qquad b\ne0.
\end{equation}
Here $b\ne0$ is precisely smoothness at $(0,0)$: it is the only
possible affine singular point since the $y$-derivative is $x$.
The standard degree-five model is smooth at infinity. Conversely
the triangular Cartier rule gives stable rank one for every $b\ne0$.

In the explicit normal form of
\citet[Proposition~2.1]{RainsSam2018Trivectors}, take
$c_9=1,c_{12}=a,c_{24}=b$, with all other coefficients zero. The tensor is
\[
 \gamma=[267]+[258]+[348]+[169]+[357]+[249]+[178]+[456]
                         +[148]+a[147]+b[134].
\]
Propositions~2.13 and 3.8 of that source identify its rank-four
surface and Coble cubic with the Jacobian of
\eqref{eq:p-rank-one-normal-form-candidate}, and show $\gamma$ is stable.
The alternating polar of its Pfaffian cubic is
\begin{align*}
 B={}&[158]+[167]+[179]+[247]+[259]+[268]+[349]+[367]+[457]\\
    &+a\bigl([148]+[157]\bigr)+b\bigl([126]+[135]\bigr).
\end{align*}
Form the infinitesimal action matrix
$\mathcal D_B:\operatorname{End}(k^9)\to\bigwedge^3k^9$ as in
Lemma~\ref{lem:p-rank-zero-genus-two-polar-stable-candidate}.
Deleting rows $[123],[345],[356]$ gives the exact polynomial identity
\begin{equation}\label{eq:p-rank-one-polar-minor-candidate}
                         \det\mathcal D'_B=b^4
                        \quad\text{in }\mathbf F_2[a,b].
\end{equation}
The certificate below specifies and checks the complete computation.
Since $b\ne0$, the Lie stabilizer of $B$ is zero. As in the preceding
lemma, stability of $\gamma$ and Corollary~5.6 of the cited source
place $B$, up to change of basis and scalar, in the Cartan subspace.
Propositions~2.9--2.10 then turn zero Lie stabilizer into stability:
every nonstable Cartan tensor has a positive-dimensional stabilizer.
\end{proof}

\begin{remark}[The uniform nonzero minor]
\label{rem:p-rank-one-polar-unit-certificate-candidate}
The complete exact program and output are
\path{agents/monodromy/code/ppav_two_rank_one_polar_uniform_minor.py}
and
\path{agents/monodromy/outputs/ppav_two_rank_one_polar_uniform_minor.json}.
They independently build all nine Pfaffians from $\gamma$, extract
$B$, and build its infinitesimal action by the exterior-product rule.
Over $\mathbf F_2[a,b]$, 77 unit pivots leave eight rows and four
columns. Four of those rows, in their retained order, have the form
\[
 \begin{pmatrix}
 b&ab&A&b^2\\
 b&ab+1&A&b^2\\
 ab+1&a^2b&D&ab^2\\
 0&0&b^2&0
 \end{pmatrix},
\quad
 \begin{aligned}
 A&=a^2b^4+b^5,\\
 D&=a^3b^4+a^2b^3+ab^5+ab^2+b.
 \end{aligned}
\]
Adding its first two rows and expanding gives determinant $b^4$.
The 77 pivot rows together with these four rows are precisely all
exterior rows except $[123],[345],[356]$; the added trace row is
unused. Hence this is a certificate of
\eqref{eq:p-rank-one-polar-minor-candidate} over the polynomial ring,
without specialization. The successful author run at
18:36:54.364729 UTC on 19 November 2026 took
0.01662725299684098 seconds. A first execution failed before any
calculation because a generated assertion omitted a space;
that exact failed source and diagnostic are retained in the history
and outputs directories. The earlier exact exploratory residual-minor
calculation and its original executed source/output are also retained.
No finite sampling is used in the proof. The parent, reduction/descent,
and obstruction/crystal agents independently audited the complete proof
and original source passages. Independent polynomial reconstructions of
the displayed tensor and fixed minor took respectively
0.04347756900097011, 0.08703309600241482, and
0.0531565969940857 seconds; reduction/descent also independently
reconstructed all nine Pfaffians. Their author-output-preserving
executions of the complete retained program took
0.019327093999891076, 0.019681259997014422, and
0.019665626001369674 seconds. All checks passed.
\end{remark}

\begin{theorem}[Every fixed nonordinary principal surface in degree one]
\label{thm:every-nonordinary-ppav-degree-one-full-candidate}
Let $(A,\Theta)$ be any nonordinary principally polarized abelian
surface over a perfect field $k$ of characteristic two.
For the smooth maximal-rank locus in $|3\Theta|$, the relative
unit local system has rank eight and both its geometric and
arithmetic images are $\GL_8(\mathbf F_2)$.
\end{theorem}
\begin{proof}
The geometric two-rank-zero case is
Theorem~\ref{thm:p-rank-zero-ppav-degree-one-full-candidate}.
If the geometric two-rank is one and the principally polarized
surface decomposes geometrically, it is the product of an ordinary
and a supersingular elliptic curve, so
Theorem~\ref{thm:mixed-every-fixed-modulus-degree-one-candidate}
applies after identifying the principal polarization factors.

Otherwise $A^{(2)}$ is the Jacobian of its smooth genus-two theta
curve, whose two-rank is one. The preceding lemma proves stability
of the polar tensor of its Coble cubic. The exact Cartier--Coble
identification in
Lemma~\ref{lem:root-nonordinary-Coble-identification-candidate}
transfers this statement to the actual Cartier Pfaffian cubic under
the complete three-theta embedding. Taking coefficient square roots
preserves stability. Thus all hypotheses of
Theorem~\ref{thm:nonordinary-ppav-polar-stable-conditional-candidate}
are satisfied, and it gives the full geometric image. The arithmetic
image is then the same full general linear group. All identifications
of principal line-bundle representatives are by translations on the
geometric abelian surface, so they identify the complete linear-system
relative local systems as well.
\end{proof}

\begin{remark}[The next fixed-surface gate]
\label{rem:nonordinary-ppav-degree-one-next-gate-candidate}
This closes degree one on every fixed nonordinary principally
polarized surface. Higher-degree structural images on a fixed
nonproduct nonordinary surface remain separate. In particular this
result does not extend the product all-degree theorems to arbitrary
principal surfaces by specialization or by padding a nodal norm
with a positive-defect ample component.
\end{remark}
